\documentclass[10pt,a4paper]{amsart}
\usepackage[margin=19mm]{geometry}
\usepackage{amsmath,amssymb,mathtools}
\usepackage[scr=rsfs,cal=euler]{mathalfa}
\usepackage{xfrac}
\usepackage[unicode,hidelinks,bookmarks=false]{hyperref}
\newcommand{\ie}{i.e.\ }
\newcommand{\cf}{cf.\ }
\newcommand{\resp}{respectively\ }
\newcommand{\Subsection}{\subsection}
\providecommand{\nobreakdash}{\nobreak\hbox{-}}
\setlength{\emergencystretch}{2em}
\multlinegap0pt
\usepackage[shortlabels]{enumitem}
\usepackage{amssymb}
\usepackage{mathtools}
\usepackage{algorithm}\usepackage{algpseudocode}
\newtheorem{lemma}{Lemma}
\newtheorem{theorem}{Theorem}
\newcommand{\HD}[2]{H^{#1}\mathbf D(#2)} % shorthand for H^k D(A)
\title{Ethic Duality: A Homological Framework for Primal-Dual Problems}
\author{Dmitry Pasechnyuk-Vilensky, Martin Tak\'a\v{c}}
\date{Source: arXiv:2512.17170v1 (CC BY 4.0)}
\begin{document}
\maketitle

\noindent\emph{Source and licence.} Ethic Duality: A Homological Framework for Primal-Dual Problems. Version arXiv:2512.17170v1 (CC BY 4.0), distributed by arXiv under the Creative Commons Attribution 4.0 International licence, http://creativecommons.org/licenses/by/4.0/, as recorded in the arXiv licence metadata for this exact version and shown on the abstract page https://arxiv.org/abs/2512.17170v1. Full licence notice: \texttt{licenses/arxiv-2512.17170-CC-BY-4.0.txt}.\par\smallskip

\noindent\textit{Editorial scope.} Counted unit: the article's theorem theorem (source0.tex lines 2025--2027). The article's complete original proof of it is delivered with it (source0.tex lines 2029--2056, one \texttt{proof} environment of 28 lines). No mathematics is written, repaired or re-derived: the statement and the proof are the article's own lines, in the article's own notation, and no retained line is substituted or re-flowed.  Retained uncounted as the source-local context the proof needs: lines 1965--1973; lines 1986--1988; lines 1994--1996; lines 2007--2009; lines 2017--2019.  Nothing was omitted from any retained span, so the package carries no omission record.  No retained line contains a citation, so the wrapper carries no bibliography and no bibliography entry.  No retained line contains a figure environment, an image-inclusion command or drawing code, so no image is delivered and the package declares no asset dependency.  Editorial formatting only, in addition: an AMS \texttt{amsart} preamble that restates the article's own declarations and macros used by the retained lines, namely the environments \texttt{lemma}, \texttt{theorem} on separate counters, together with the standard packages those lines need.  The retained environments print in this document as Lemma 1, Theorem 1 in order of appearance, whereas the article prints the same items with the numbering of its own full document.  The complete native source of the article as distributed in the arXiv e-print for this exact version is bundled unchanged as \texttt{sources/191296/source0.tex}.
\begin{lemma}[Verifier capture via Cook--Levin]\label{lem:verifier-capture}
Let $L\in \mathbf{NP}_{\mathrm{TM}}$. Then there exist polynomials $q_{\mathrm{ver}}(n)$ and $p(n)$ such that, for the bar/cobar model $D_{L,n}^{\bullet}$ of §3.3, the truncation depth $q_{\mathrm{ver}}(n)$ satisfies
\[
H^{0}\mathbf D\!\big((D_{L,n}^{\bullet})_{\le q_{\mathrm{ver}}(n)}\big)(x,w)=1
\iff
\text{$w\in\{0,1\}^{\le p(n)}$ is an accepting witness for $x\in\{0,1\}^n$}.
\]
Moreover, the map $1^n\mapsto \langle (D_{L,n}^{\bullet})_{\le q_{\mathrm{ver}}(n)}\rangle$ is logspace–uniform.
\end{lemma}

\begin{lemma}[Uniform depth–time bridge for \(\mathbf P\)]\label{lem:PTM-circuits}
A language \(L\) lies in \(\mathbf P_{\mathrm{TM}}\) iff there exists a logspace–uniform family of Boolean circuits \(\{C_n\}\) of size \(n^{O(1)}\) and depth \(n^{O(1)}\) computing \(f_{L}\) on \(\{0,1\}^{n}\).
\end{lemma}

\begin{lemma}[Ethic verifiability implies Turing verifiability]\label{lem:EthNP-to-NPTM}
If \(L\in \mathbf{NP}_{\eth}\), then \(L\in \mathbf{NP}_{\mathrm{TM}}\).
\end{lemma}

\begin{lemma}[Ethic polynomial reconstruction compiles to poly–size, poly–depth circuits]\label{lem:EthP-to-circuits}
If \(L\in \mathbf P_{\eth}\) with witness \(Q\), then there exists a logspace–uniform family \(\{C_n\}\) of circuits of size \(n^{O(1)}\) and depth \(O(Q(n))\) computing \(f_{L,n}\).
\end{lemma}

\begin{lemma}[Ethic depth equals formula/KW depth]\label{lem:KW}
For each \(n\), the least \(q\) with \(H^{>q}\mathbf D(D_{L,n}^{\bullet})=0\) equals the minimum Boolean formula depth and the deterministic Karchmer–Wigderson depth of \(f_{L,n}\).
\end{lemma}

\begin{theorem}[Equivalence of ethic and Turing \(P{=}NP\)]
\(\mathbf P_{\eth}=\mathbf{NP}_{\eth}\) if and only if \(\mathbf P_{\mathrm{TM}}=\mathbf{NP}_{\mathrm{TM}}\).
\end{theorem}

\begin{proof}
(\(\Rightarrow\)).
Assume \(\mathbf P_{\eth}=\mathbf{NP}_{\eth}\).
Let \(L\in \mathbf{NP}_{\mathrm{TM}}\) be arbitrary.
By Lemma~\ref{lem:verifier-capture} (Verifier capture via Cook--Levin), there exist polynomials \(q_{\mathrm{ver}}(n)\) and \(p(n)\) such that, for every \(n\),
\[
x\in L \;\Longleftrightarrow\; \exists\,w\in\{0,1\}^{\le p(n)}:\ 
H^{0}\mathbf D\!\big((D_{L,n}^{\bullet})_{\le q_{\mathrm{ver}}(n)}\big)(x,w)=1 .
\]
By \emph{(M1)}–\emph{(M2)} (Verification extraction), the truncations \((D_{L,n}^{\bullet})_{\le q_{\mathrm{ver}}(n)}\) compile into a logspace–uniform family of verifier circuits \(V_n\) of size \(\mathrm{poly}(n)\) with
\[
x\in L \;\Longleftrightarrow\; \exists\,w\in\{0,1\}^{\le p(n)}:\ V_n(x,w)=1 .
\]
Hence \(L\in \mathbf{NP}_{\eth}\) by the definition of \(\mathbf{NP}_{\eth}\).
By the hypothesis \(\mathbf P_{\eth}=\mathbf{NP}_{\eth}\), we conclude \(L\in \mathbf P_{\eth}\).
Therefore there exists a polynomial \(Q(n)\) such that \(\HD{1}{(D_{L,n}^{\bullet})_{\le Q(n)}}=0\) for all \(n\) and the computation extractor \emph{(M4)} yields circuits \(C_n:=C_{n,Q(n)}\) computing \(f_{L,n}\) of depth \(O(Q(n))\) and size \(n^{O(1)}\), logspace–uniform in \(n\).
By Lemma~\ref{lem:EthP-to-circuits}, the family \(\{C_n\}\) is poly–size and poly–depth, and by the uniform circuit characterization of \(\mathbf P\) (Lemma~\ref{lem:PTM-circuits}) we obtain \(L\in \mathbf P_{\mathrm{TM}}\).
Since \(L\in \mathbf{NP}_{\mathrm{TM}}\) was arbitrary, we have \(\mathbf{NP}_{\mathrm{TM}}\subseteq \mathbf P_{\mathrm{TM}}\), hence \(\mathbf P_{\mathrm{TM}}=\mathbf{NP}_{\mathrm{TM}}\).

(\(\Leftarrow\)).
Assume \(\mathbf P_{\mathrm{TM}}=\mathbf{NP}_{\mathrm{TM}}\).
Let \(L\in \mathbf{NP}_{\eth}\).
By Lemma~\ref{lem:EthNP-to-NPTM}, \(L\in \mathbf{NP}_{\mathrm{TM}}\), hence \(L\in \mathbf P_{\mathrm{TM}}\).
By Lemma~\ref{lem:PTM-circuits} there exists a logspace–uniform poly–size, poly–depth circuit family \(\{C_n\}\) computing \(f_{L}\).
By Lemma~\ref{lem:KW} the minimal ethic internal depth equals the formula/KW depth and is polynomially bounded in \(n\); therefore there is a polynomial \(Q\) such that \(H^{>Q(n)}\mathbf D(D_{L,n}^{\bullet})=0\) for all \(n\), whence \(\Delta_{L}(n,Q(n))=0\).
Finally, \emph{(M4)} produces the required logspace–uniform circuits \(C_{n,Q(n)}\) and the definition of \(\mathbf P_{\eth}\) is met.
Thus \(\mathbf{NP}_{\eth}\subseteq \mathbf P_{\eth}\) and equality follows.
\end{proof}

\end{document}
